\section{Rules checked through the public record} \label{app:public}

The scripts check the software environment, the sector, the star, the file name, the registered arguments and the
checksums that every result carries (Appendices~\ref{app:localisation} and \ref{app:procedure}). They do \emph{not}
check the rules of this section. These rules are part of the pre-registration exactly like the coded ones. They are
checked through the public record: every file named below is posted with its sha256 as in
Appendix~\ref{app:procedure}, point 3, and because the data and the code are public, anyone can repeat each step.
Nothing in this section claims that the code prevents a violation.

\subsection{Invalid results} \label{app:p1}

\begin{itemize}
\item A result (a photometric, variant or localisation output file, a summary file, or a sentence) obtained in
  violation of a pre-registered rule on the ephemeris and the predicted times (Sections~\ref{sec:ephem} and
  \ref{sec:times}), the statistic (Section~\ref{sec:stat}), or of a rule of Appendices~\ref{app:data}, \ref{app:processing}, \ref{app:variants},
  \ref{app:localisation}, \ref{app:procedure}, \ref{app:public} or \ref{app:commands}, is \emph{invalid}. An invalid result is not an
  outcome of this test. It cannot be used to accept or to reject signal A or signal B, and it is not described with the
  sentences of Appendix~\ref{app:wording}. It is still posted, labelled ``invalid'' with the rule it breaks
  (Appendix~\ref{app:procedure}, point 5).
\item That a script accepted a run, or that \code{s110\_summary.py} wrote a sentence, does not by itself make the result
  valid.
\item If a violation is found after posting, by me or by anyone, a new version of this record declares the affected
  results invalid and names the rule. Where the rules still allow it, the step is repeated as registered (for example
  the primary run on the product that Appendix~\ref{app:data} requires). An invalid run never counts as the first run of
  Appendix~\ref{app:procedure}, point 3; the repeated run defines the outcome, and the delay and the earlier invalid run
  are reported. Where the rules no longer allow a repeat, the outcome of each affected signal is O4 with the reason
  ``pre-registered procedure not followed'', or, for the localisation, ``localisation not possible'' with that reason.
  This is the case when the date 2027 February 28 or a reading limit of Appendix~\ref{app:p5} has passed, or when a rule
  of order of Appendix~\ref{app:procedure}, points 1 and 3, was broken, which a repeat cannot repair.
\item The time limits of Appendix~\ref{app:procedure}, points 1 and 3, are administrative: the primary run within 30 days
  of the release of the primary product; results posted within 7 days; variants and localisation run within 14 days of
  posting the primary result. No outcome, threshold, label or sentence depends on them, and none of the scripts reads
  them (the only date any script compares is 2027 February 28, in \code{s110\_summary.py}). A run or a posting after
  such a limit is reported with its delay; it is not invalid for that reason alone. Binding, as before, are every rule
  of order (no Sector 110 product opened before the primary run, Appendix~\ref{app:procedure}, point 1; variants and
  localisation only after the primary result is posted, point 3), the date 2027 February 28
  (Appendices~\ref{app:data}, \ref{app:variants}, \ref{app:p2} and \ref{app:p5}), and the reading limits of
  Appendix~\ref{app:p5}, which decide an outcome.
\item Two different 30-day limits exist, and they are never interchangeable.
  \begin{itemize}
  \item The \emph{run limit} (Appendix~\ref{app:procedure}, point 1): the primary run is made within 30 days of the
    release of the primary product. It is administrative, as stated above: a later primary run is still the primary
    run, with its delay reported.
  \item The \emph{reading limits} (Appendix~\ref{app:p5}): if a file cannot be read by the registered script, an adapted
    reader that passes the check of Appendix~\ref{app:procedure}, point 4, must be posted within 30 days of the release
    of that file at MAST. The file is the light curve for the photometric run and the target pixel file for the
    localisation. The release date is the one listed by the product search record of Appendix~\ref{app:p2}, or, where
    MAST lists none, the date of the first product search that lists the file. The limit ends at 24:00 UTC of the 30th
    day. The reading limits are binding: when one has passed without such a posting, the outcome of the corresponding
    row of Table~\ref{tab:failures} applies.
  \end{itemize}
  Meeting the run limit never meets a reading limit, and the administrative status of the run limit never applies to a
  reading limit. Each of the two reading limits counts only from the release of its own file.
\end{itemize}

\subsection{Choice of the product (fallbacks)} \label{app:p2}

\begin{itemize}
\item Before the primary run a product search record is saved. It contains the exact MAST query (the astroquery call
  with all its arguments, or the MAST Portal search with its URL), the UTC date and time of the search, and the
  complete, unfiltered list of Sector 110 products of TIC 4206066 that the search returns (file names, product types,
  and the release date or data release number where MAST gives one). It names the light curve and the target pixel
  file selected, with their sha256 checksums, and, for a renamed file (Appendix~\ref{app:procedure}, point 2), the
  original name. Listing, downloading and computing checksums of these files is not an opening of the product in the
  sense of Appendix~\ref{app:procedure}, point 1.
\item The primary run is valid only on the first released SPOC 2-min light curve listed by this record. A fallback run
  is valid only if the record was made on or after 2027 February 28 and lists no SPOC 2-min light curve (fallback 1,
  QLP), or lists neither a SPOC 2-min nor a QLP light curve (fallback 2, TESS-SPOC). A fallback run made before
  2027 February 28, without this record, or while the record lists a product of higher priority is invalid, whatever
  its outcome.
\item The code does not check any of this: \code{s110\_final.py} accepts a QLP or TESS-SPOC file with a registered file
  name as a registered run on any date, and \code{s110\_summary.py} writes sentences for it.
\item Variant (d) uses a record of the same kind for its QLP file (Appendix~\ref{app:variants}). The record for the
  primary product is posted with the primary result. The record for (d) is posted with the result of (d), or, if (d) is
  not run, with the summary that declares it not available.
\end{itemize}

\subsection{Time system of every input file} \label{app:p3}

\begin{itemize}
\item Every analysed file must declare in its first extension (HDU 1) the time system of all archival inputs of this
  test, which was checked on the archival SPOC 2-min, TESS-SPOC and QLP files of Sectors 6, 32 and 98 on 2026 October 5:
  \code{TIMESYS} $=$ TDB, \code{TIMEREF} $=$ SOLARSYSTEM, \code{BJDREFI} $=$ 2457000, \code{BJDREFF} $=$ 0 or absent (QLP
  files do not carry it), \code{TIMEUNIT} $=$ d, and \code{TIMEDEL} equal, within 1\,\%, to 2\,min for a SPOC 2-min light
  curve or target pixel file and to 200\,s for a QLP or TESS-SPOC 200-s light curve.
\item As the first step of every run, before the registered command, these keys are printed without reading any data
  column, with the command given in Appendix~\ref{app:commands}. The printed line is kept in the attempt record
  (Appendix~\ref{app:p4}). Reading these keys is part of the run, not an opening of the product in the sense of
  Appendix~\ref{app:procedure}, point 1.
\item If any key differs, the file is not analysed, no conversion is made, and there is no switch to another product.
  For the primary file the outcome is O4 for both signals, with the reason ``time system of the product not as
  registered''. For a variant file the variant is not available, with the printed line as its record. For the pixel
  file the localisation is not possible, with the reason ``time system of the pixel file not as registered''.
\item A run on a file whose keys were not printed and kept, or differ, is invalid. \code{s110\_final.py} and
  \code{s110\_localise.py} do not read these keys.
\end{itemize}

\subsection{Attempt records and ``not available''} \label{app:p4}

\begin{itemize}
\item Every attempt of every registered command (primary or fallback, each variant, each localisation, the summary)
  has a plain-text attempt record, written while the run happens and not edited afterwards. It contains the exact
  command line; the name and sha256 of every input file; the sha256 of the script run; for a run on a FITS file, the
  printed line of Appendix~\ref{app:p3}; the UTC start and end time; the exit status; and the complete standard output and
  standard error. The record is written by the macOS \code{script} utility in UTC, for example
  \code{TZ=UTC script -a RECORD <command>}, which saves the command, its start and end time, all its output and its exit
  status. The checksum lines of the inputs and of the script, and the line of Appendix~\ref{app:p3}, are written into
  the same record in the same way just before the registered command. There is one record per attempt, also for
  successful attempts and for refusals. All records are posted with the output files they belong to
  (Appendix~\ref{app:procedure}, point 3), so the number and the order of the attempts are public.
\item The variants are attempted in one continuous session, in the order (a), (b), (c), (e), (f), and then (d) if a QLP
  Sector 110 release is public when the session starts. Each is attempted once and, after a failure, repeated once
  (Appendix~\ref{app:procedure}, point 4). The order is binding. A session that runs past midnight UTC continues in the
  same order; its start and end are in the records, and nothing else follows from the date. If no QLP Sector 110
  release is public when the session starts, (d) is attempted later, in the same way and with its own records, on the
  QLP release that is public on 2027 February 28 (Appendix~\ref{app:variants}), on or after that date. If MAST lists no
  QLP Sector 110 release on 2027 February 28, (d) is final as ``not-public-on-deadline''.
\item A variant is ``not available'' only if its two records show the registered command on the registered input (for
  (a) to (c), (e) and (f), the primary file, with the same sha256), each ending with a non-zero exit status after an
  error while reading or computing. For a reading failure, which Appendix~\ref{app:procedure}, point 4, does not repeat,
  one such record is enough. A variant is also not available if its file fails the check of Appendix~\ref{app:p3}. The
  file given to \code{--not-available} is the record of the last attempt (or the printed line of
  Appendix~\ref{app:p3}).
\item A variant whose registered run wrote its output file can never be declared not available; that file is its
  result. A refusal is not ``not available'' (Appendix~\ref{app:variants}).
\item A summary that declares a variant not available with any other file (free text, a refusal, an edited or
  shortened output, the record of another command) is invalid. \code{s110\_summary.py} accepts any non-empty file there
  and does not check its content.
\item Variant (d) is final as ``crashed'' with the two records of the attempt on a QLP release that was public at the
  time. It is final as ``not-public-on-deadline'' with the product search record of Appendix~\ref{app:p2}, made on or
  after 2027 February 28, that lists no QLP Sector 110 light curve. A file for (d) that fails the check of
  Appendix~\ref{app:p3} is declared final with \code{--d-final-reason crashed} and the printed line of
  Appendix~\ref{app:p3} as the file.
\item When a variant among (a) to (e) is not available, the robustness qualifier is computed without it. The line of
  Appendix~\ref{app:p7} shows it as ``not available'', and the qualifier is never described as covering it.
\end{itemize}

\subsection{Registered failures that leave no output file} \label{app:p5}

Some registered outcomes are reached without an output file, so \code{s110\_summary.py} cannot write their sentences
(it needs the primary result and, for a SPOC primary, the localisation result of each signal with O1 or O2). For these
cases only, the sentence is written from the posted files exactly by the tables of Appendix~\ref{app:wording}; only the
reason comes from Table~\ref{tab:failures}.

\begin{longtable}{>{\raggedright\arraybackslash}p{0.34\textwidth}>{\raggedright\arraybackslash}p{0.16\textwidth}>{\raggedright\arraybackslash}p{0.14\textwidth}>{\raggedright\arraybackslash}p{0.26\textwidth}}
\caption{Registered failures that leave no output file.\label{tab:failures}}\\
\toprule
Case & Posted record & Outcome & Fixed reason R \\
\midrule
\endfirsthead
\toprule
Case & Posted record & Outcome & Fixed reason R \\
\midrule
\endhead
the primary photometric run ends twice without writing its output (not a refusal, not a reading failure;
  Appendix~\ref{app:procedure}, point 4) & the two attempt records & O4 for A and B & ``the registered photometric run
  failed twice without writing a result'' \\
the primary file cannot be read, and no adapted reader (Appendix~\ref{app:p6}) that passes the check of
  Appendix~\ref{app:procedure}, point 4, is posted within 30 days of the release of the primary product & the attempt
  record with the error output & O4 for A and B & ``the registered product could not be read'' \\
the primary file fails the check of Appendix~\ref{app:p3} & the printed line & O4 for A and B & ``time system of the
  product not as registered'' \\
localisation is needed, and a product search of 2027 February 28 or later lists no SPOC 2-min target pixel file of
  Sector 110 & that search record & localisation not possible & ``no SPOC 2-min target pixel file released by
  2027-02-28'' \\
the localisation run ends twice without writing its output & the two attempt records & localisation not possible &
  ``the registered localisation run failed twice without writing a result'' \\
the pixel file cannot be read, and no adapted reader is posted within 30 days of its release & the attempt record
  with the error output & localisation not possible & ``the target pixel file could not be read'' \\
the pixel file fails the check of Appendix~\ref{app:p3} & the printed line & localisation not possible & ``time system
  of the pixel file not as registered'' \\
\bottomrule
\end{longtable}

\begin{itemize}
\item Photometric failure: for each signal X, ``Sector 110 cannot decide on signal X (reason: R).''
\item Localisation failure: the photometric part and the robustness part are written exactly as
  Appendix~\ref{app:wording} defines them, from the posted primary and variant results and
  \code{calib/false\_alarm\_levels.json}. The localisation part is ``; its position could not be measured (R).'' for A
  and ``; B's localisation was not possible (R).'' for B, written like the localisation parts of
  Appendix~\ref{app:wording}, without the word ``reason''. The sentence of the other signal is written in the same way
  from its own files.
\item These sentences are not produced or checked by the code; anyone can rebuild them from the tables and the posted
  files.
\item A pixel file that is missing before 2027 February 28 is not a final failure: the sentences wait until the file is
  released or the date passes. The primary result is posted in any case (Appendix~\ref{app:procedure}, point 3).
\item A failure never replaces a result: if any registered attempt of a step wrote its output file, that file defines
  the result (Appendix~\ref{app:procedure}, point 4), and a failure sentence for that step is invalid.
\end{itemize}

\subsection{A reader adapted under the failure rule} \label{app:p6}

\begin{itemize}
\item \code{s110\_summary.py} compares the checksums carried by each result with the scripts in the folder where it
  runs, under their registered names \code{s110\_final.py} and \code{s110\_localise.py}. An adapted script therefore
  keeps its registered file name and replaces the original in the copy of this deposit that is used for the Sector 110
  runs.
\item The adapted photometric script is used unchanged for the primary and for every variant of that product. This is
  not an edit of a variant: variants are never edited individually.
\item The helper modules in \code{deps/} are never adapted. A pixel file that \code{deps/pixloc\_lib.py} cannot read is
  the case ``the target pixel file could not be read'' of Table~\ref{tab:failures}.
\item The sha256 of the adapted script, its diff against the deposited script, and the check of
  Appendix~\ref{app:procedure}, point 4, are posted before the Sector 110 run. A result whose \code{script\_sha256} is
  neither the checksum in \code{SHA256SUMS} nor a checksum posted in this way is invalid.
\item The summary compares only with the file on its own disk, so it accepts any script placed there; the check that
  counts is the public comparison of Appendix~\ref{app:procedure}, point 7.
\end{itemize}

\subsection{What ``robust'' means, and the line of variant outcomes} \label{app:p7}

\begin{itemize}
\item In the sentences, ``robust'' means exactly this: among the variants (a) to (e) that were run and reached
  $\Sens \geq 0.95$, none returned the opposite outcome (O3 against a primary O1; O1 or O2 against a primary O3). It
  does \emph{not} mean that all variants agree. A qualifying variant that returns O2 (a dip at the wrong depth or time)
  or O4 against a primary O1, or O4 against a primary O3, leaves ``robust'' in place. Variants below 0.95, variant (f),
  and variants that are not available or pending are not part of it.
\item Every posted summary is followed, for each signal, by one fixed line written by hand from the posted variant
  results and the summary file (the script does not write it). An example: ``Variant outcomes for signal A: a O1 (0.97);
  b O2, timing off (0.99); c O1 (0.96); d pending; e not available; f O1 (0.95; comparison only).'' The line lists every
  applicable variant in the order a to f. For each it gives the label O1, O2, O3 or O4 taken from the start of the
  outcome string of the script (for O2 followed by the failed conditions written by the script, for O4 by its reason)
  and its sensitivity to two decimals where its output file has one, else ``no sensitivity computed''; or ``not
  available''; or ``pending''. Variant (f) is listed in the same way, marked ``comparison only''. Wherever a sentence
  containing ``robust'' is quoted, this line is quoted with it.
\item ``Robust'' is never rewritten as ``confirmed by all processings'' or in any wider sense.
\end{itemize}
